\documentclass[../main.tex]{subfiles}
\begin{document}
%==========================================TIÊU ĐỀ TẬP========================================
\begin{table}[h]
  \begin{adjustwidth}{-1cm}{}
    \begin{tabular}{>{\centering\arraybackslash}p{6cm}|>{\raggedright\arraybackslash}p{14cm}}
      \multirow{5}{6cm}
      % Cột bên trái: ÉPISODE và số tập
      {\\[-40pt]
        \flaregothic\fontsize{20pt}{20pt}\selectfont \centering
        \color{eptype}ÉPISODE\color{black}\\
        \burbank\fontsize{72pt}{72pt}\selectfont
        \color{epnum}5\color{black}\\[6pt]
        \cafeteria\fontsize{20pt}{20pt}\selectfont\color{epnum}(1-1)\color{black}
        \\[18pt]
        \flaregothic\fontsize{14pt}{14pt}\selectfont
        \color{epattr}BIENVENUE, \uline{\color{Banpire} Mel} !
        \flaregothic\fontsize{18pt}{18pt}\selectfont
        \color{epattr}ÉPISODE PERDU
        \color{black}
      }
       &
      % Dòng thứ nhất: tên tập tiếng Nhật
      {\fotskip\fontsize{18pt}{18pt}\selectfont \ruby{右}{みぎ}\ruby{腕}{うで}が\ruby{意}{い}\ruby{思}{し}を\ruby{持}{も}って\ruby{走}{はし}り\ruby{出}{だ}す【\ruby{悲}{ひ}\ruby{報}{ほう}】
      } \\[6pt]
      % Dòng thứ hai: tên tập tiếng Anh
       & {\cafeteria\fontsize{18pt}{18pt}\selectfont In VR, ARM Waves YOU}                                                                                     \\[4pt]
      % Dòng thứ ba: tên tập tiếng Việt
       & {\vnmsans\fontsize{18pt}{18pt}\selectfont Cánh tay bị rời của Roboco}                                                                                 \\[8pt]
      % Dòng thứ tư: Nhân vật xuất hiện
       & {\caviar{\fontsize{14pt}{14pt}\selectfont Nhân vật: \emp{kuon1} ({\abv Truyện phụ}) \emp{roboco} \emp{mel} \emp{fubuki} \emp{haato} \emp{matsuri}}}    \\[8pt]
      % Dòng cuối cùng: Thời gian diễn ra của tập (thường sẽ là ngày phát hành)
       & {\continuum\fontsize{16pt}{16pt}\selectfont Ngày phát hành: 02/06/2019}                                                                               \\[18pt]
    \end{tabular}
  \end{adjustwidth}
\end{table}
\color{black}

\txtbx[2]{
  {\color{vol} \uline{Tóm tắt:}} Một bưu kiện bí ẩn không ghi người gửi được chuyển đến văn phòng \hll\ khiến các cô gái vô cùng tò mò. Trong lúc giúp đỡ mọi người nâng chiếc hộp, cô nàng rô-bốt Roboco vô tình làm rơi mất cánh tay phải, khơi mào cho một cuộc rượt đuổi náo loạn khắp văn phòng. Kỳ lạ thay, bên trong chiếc hộp lại chính là cánh tay thay thế hoàn hảo được gửi kèm lá thư chu đáo từ người quản lý Hikawa Kuon. Trong khi đó, cánh tay cũ tự động bò về xưởng cơ khí nhà Hikawa và được hai bố con Kuon tiễn đưa dưới lò dung nham nóng chảy như một thước phim đầy cảm xúc.
}

Những ngày đầu hè đã gõ cửa thành phố Akiba ảo, mang theo cái nắng chói chang và luồng không khí oi bức đặc trưng tràn vào từng ngóc ngách của văn phòng \hll\ mới mở.

Do hệ thống điều hòa không khí vẫn chưa được lắp đặt hoàn thiện, thứ duy nhất cứu rỗi các cô gái khỏi cái nóng hầm hập lúc này chỉ là vài chiếc quạt máy đang chạy hết công suất. Tiếng cánh quạt vù vù hòa lẫn với tiếng trò chuyện rôm rả tạo nên một bầu không khí ồn ào nhưng lại mang đậm sức sống của tuổi trẻ.

Fubuki, với đôi tai cáo trắng thỉnh thoảng lại vểnh lên nghe ngóng, đang ngồi cặm cụi trước màn hình máy tính. Cô nàng đang xem xét lại bản danh sách dài thườn thượt những món đồ dùng cần mua sắm cho văn phòng mới, trong đó có cả một vài ``yêu cầu đặc biệt'' từ một tiền bối tên R.

Vừa gõ phím lạch cạch, nàng Cáo trắng vừa thỏa mãn liếm láp cây kem vị dưa gang mát lạnh, thỉnh thoảng lại rướn người về phía trước để hứng trọn luồng gió hiếm hoi từ chiếc quạt bàn nhỏ xíu.

*KÍNH COONG*

Tiếng chuông cửa bất ngờ vang lên, cắt đứt dòng suy nghĩ của nàng Cáo.

``Đến ngay đây\~{}!'' cô nàng vội vàng đáp lời, lật đật chạy ra sảnh.

Sau vài thủ tục xác nhận đơn giản, nhân viên giao hàng trao cho cô một thùng các-tông lớn, dài cỡ bằng sải tay người lớn.

``Đây, chữ ký của em đây! Cảm ơn anh nha\~{}''

Vừa dứt lời, Fubuki vô tư dùng hai tay nhấc bổng chiếc hộp lên. Nhưng cô đã đánh giá quá thấp trọng lượng của nó. Khối lượng khổng lồ bất ngờ trì kéo xuống khiến nàng Cáo mất đà, loạng choạng lùi lại phía sau.

Và rồi, điều tồi tệ nhất cũng xảy ra.

*BỊCH!*

``Ái\dots!''

Chiếc hộp tuột khỏi tay, rơi thẳng xuống và đè trúng ngay bàn chân trái của Fubuki. Cô nàng lập tức ngã bệt xuống sàn, hai tay ôm lấy bàn chân sưng tấy, đôi tai cáo cụp rủ rượi, nước mắt lưng tròng, còn hai má thì đỏ bừng lên vì đau đớn.

\img{1-1a}

\ct

Chắc hẳn lúc này các bạn sẽ tự hỏi: ``Vậy Hikawa Kuon đang ở đâu mà lại để một cô gái mảnh mai như Fubuki-chan phải gồng gánh một kiện hàng nặng trịch thế này?''

Câu trả lời rất đơn giản: Cậu ta không có mặt ở văn phòng.

Hôm nay, gia đình Kuon có việc đột xuất cần giải quyết, và cậu đã xin phép Giám đốc Tanigo nghỉ một hôm để về phụ giúp. Trước khi rời đi, Kuon đã cẩn thận nhờ cô nàng Cáo trắng để mắt tới văn phòng, dù sâu thẳm trong lòng, cậu thừa biết giao trọng trách này cho một nàng Cáo lắm mưu thì chẳng khác nào ``giao trứng cho ác''.

Trở lại với thực tại. Fubuki nhận ra việc nhấc bổng chiếc hộp bằng tay không là nhiệm vụ bất khả thi, nhất là khi chân trái vẫn còn đang nhức nhối. Đảo mắt quanh phòng, cô nàng quyết định đổi chiến thuật.

Cô nàng xoay người, tì sát lưng vào mặt bên của thùng các-tông, lấy thế rồi dồn toàn bộ sức lực của cơ thể để đẩy nó lết trên sàn nhà. Lúc thì dùng lưng ủi, khi thì dùng cả hai tay đẩy, nàng Cáo dốc hết sức bình sinh để di chuyển khối ``đá tảng'' này về phía chiếc bàn ở giữa phòng.

\img{1-1b}

Thế nhưng, đoạn đường chưa qua được một nửa thì Fubuki đã hoàn toàn kiệt sức. Cô nàng ngã ngửa ra sàn nhà, tứ chi dang rộng thành hình chữ đại, ngực phập phồng thở dốc từng cơn.

Chứng kiến bộ dạng thê thảm của người bạn, Natsuiro Matsuri lật đật chạy tới hỏi han.

``Fubuki, bà không sao chứ?!'' Matsuri hớt hải chạy lại, vẻ mặt đầy lo lắng.

``Tôi\dots\ tôi \dots\ ổn.'' Fubuki thều thào đáp, một tay chống hông, tay kia uể oải vẫy vẫy. ``Giá mà Kuon-senpai có ở đây lúc này thì tốt biết mấy\dots''

``Nhưng mà anh ấy lại bận việc gia đình mất tiêu rồi\dots\'' nàng Cổ động thở dài thườn thượt. ``Chán anh ấy thật chứ.''

Gần đó, Haato và Roboco đang đứng nhìn chiếc hộp các-tông lớn mà Fubuki vừa chật vật đẩy vào. Cô nàng Nữ sinh trung học khẽ nhíu mày, giọng đầy vẻ nghi hoặc: ``Thùng hàng này trông nặng thật đấy.''

Roboco ngồi xổm xuống, xem xét lớp ngoài của chiếc hộp. Cô nhận thấy chỉ có vài tem cảnh báo dán rải rác mà không hề có bất kỳ thông tin liên lạc hay địa chỉ người gửi nào. Một thùng các-tông màu nâu thô ráp đơn điệu, chỉ được niêm phong bởi một đường băng dính dán dọc theo mép.

``Trên hộp hoàn toàn không có địa chỉ người gửi hay ghi chú gì cả,'' nàng Người máy thông báo sau khi sờ soạn xung quanh.

``Vậy sao? Lạ thật đấy\dots\ Thế thì thử mở ra xem bên trong rốt cuộc là cái gì đi\dots'' Haato tò mò đề xuất.

Nhưng ngay khi Roboco bắt đầu kiểm tra chiếc hộp, một giọng nói vang lên từ phía cửa ra vào. Đó là một thành viên khác từ Thế hệ thứ Nhất, đang lon ton chạy vào và chào hỏi mọi người.

\chrint

\model{24}{l}{6cm}{mel}

Yozora Mel (\ruby{夜}{よ}\ruby{空}{ぞら}メル, \ruby{Dạ}{Yozora} Me-rư) - cô nàng Ma cà rồng đến từ Ma giới xa xôi nhưng lại mang danh hiệu ``thần đồng'' đầy bí ẩn. Điểm kỳ lạ nhất ở Mel chính là việc cô vô cùng ghét máu - thứ vốn dĩ là nguồn sống của ma cà rồng. Thay vào đó, cô nàng lại có niềm đam mê bất tận với nước ép sơ-ri (acerola) chỉ đơn giản vì chúng có màu đỏ bắt mắt. Không chỉ vậy, Mel cũng chẳng hề e ngại ánh nắng mặt trời rực rỡ như đồng loại của mình\footnote{Quyển số Một - tập 15.}!

Khác xa với hình ảnh lạnh lùng, khát máu thường thấy của loài ma cà rồng trong truyền thuyết, Mel sở hữu một tính cách ngọt ngào, điềm tĩnh và vô cùng chu đáo. Cô luôn sẵn lòng đặt lợi ích của người khác lên trên bản thân mình, dù điều đó đôi khi vô tình đẩy cô vào những rắc rối dở khóc dở cười vì tính nhường nhịn cả với những đàn em.

Nhờ sự dịu dàng thiên bẩm, Mel sở hữu chất giọng ấm áp, mềm mại và vô cùng êm tai. Tuy nhiên, tông giọng của cô nhỏ đến mức Fubuki từng trêu rằng: ``Mỗi lần nói chuyện, Kuon-senpai toàn phải ghé sát tai vào mới nghe rõ rồi lắc đầu bảo cậu ấy nói to lên.'' Phong cách nói chuyện nhỏ nhẹ này của Mel có phần tương đồng với cô em Yuzuki Choco.

Ngoại hình của cô nàng cũng vô cùng nổi bật với mái tóc bob màu vàng óng ả điểm thêm sợi tóc mái tinh nghịch. Mái tóc ấy được tô điểm bởi một chiếc kẹp hình dơi cùng viên đá quý đỏ lấp lánh, hài hòa với đôi mắt hổ phách sâu thẳm. Trang phục của Mel mang phong cách quyến rũ đặc trưng với chiếc áo nịt ngực viền ren đan chéo, phía sau là đôi cánh dơi nhỏ nhắn. Hai ống tay áo màu trắng tách rời kết hợp với dây đai đen cá tính. Cô đeo găng tay hở ngón, diện quần soóc đen ngắn được che phủ khéo léo bằng lớp váy ngoài trong suốt bồng bềnh. Đôi chân thon dài được làm nổi bật bởi đôi tất không đối xứng (bên dài đến đùi có dây đai, bên chỉ cao đến đầu gối) và đôi bốt đen buộc dây cam vàng bắt mắt với cánh dơi nhỏ chìa ra hai bên.

``Kon-kappu\~{}!'' là câu chào quen thuộc của cô với những người hâm mộ khi cô bắt đầu phát sóng trực tiếp.

\chrint

``Chào buổi sáng mọi-''

Tuy nhiên, do đang mải mê chỉnh lại tóc, nàng Dơi không để ý đến chiếc hộp lớn chắn ngay lối đi. Thế rồi cô vấp phải nó và ngã ngửa ra sau, mông đập mạnh xuống sàn.

``Bà không sao chứ?'' Haato nhanh chóng đi tới cô nàng cùng Thế hệ, khuỵu người xuống hỏi han.

\img{1-1b1}

``T-Tôi ổn\dots'' nàng Ma cà rồng run rẩy đáp, đôi mắt hơi ngấn lệ.

``Cái hộp to như vậy mà cứ để ở đây thì nguy hiểm lắm,'' Haato vừa đỡ Mel dậy vừa nhắc nhở.

``Tại sao nó lại ở đây chứ?'' Mel nghiêng đầu đầy thắc mắc.

``Tôi không biết. Fubu-chan mang nó vào đây mà,'' cô nàng Nữ sinh trung học nhún vai giải thích.

Roboco, đứng gần đó, quyết định xử lý tình huống. ``Vậy để mình đặt nó lên bàn nhé.'' Biết rằng với sức mạnh vượt trội của một rô-bốt thế hệ cao, việc này không thể làm khó được cô, Fubuki liền cười động viên: ``Quả đúng là rô-bốt thế hệ cao Roboco-senpai có thể làm được mà! Bọn em nhờ cả vào chị đấy!'' Fubuki cười rạng rỡ.

``Roboco-senpai đúng là mạnh thật!'' Matsuri cũng hào hứng phụ họa.

Roboco khẽ gật đầu, chuẩn bị thể hiện sức mạnh của mình để di chuyển chiếc hộp nặng nề. Các thành viên còn lại đứng xung quanh, ánh mắt đầy mong đợi và tò mò.

Đúng như suy đoán của mọi người, chiếc hộp lớn với cô nàng chẳng phải là vấn đề gì quá to tát khi cô có thể nâng nó lên một cách nhẹ nhàng, nhưng ngay khi vừa nhấc nó lên, cánh tay phải của cô đột ngột rơi ra cùng với chiếc hộp.

*RẦM!*

Một tiếng động lớn vang lên khi cánh tay và chiếc hộp cùng rơi xuống sàn, khiến Roboco khẽ giật mình, còn các thành viên khác thì kinh ngạc không thốt nên lời. Tất cả ánh mắt đổ dồn vào cánh tay bị rơi, và mọi người (trừ Roboco) đồng loạt hét lên thất thanh.

\begin{figure}[H]
  \centering
  \begin{subfigure}[b]{0.75\textwidth}
    \centering
    \includegraphics[width=12.5cm]{../../../../Image/Book1/1-1c1.png}
  \end{subfigure}
  \quad
  \quad
  \begin{subfigure}[b]{0.45\textwidth}
    \centering
    \includegraphics[width=8cm]{../../../../Image/Book1/1-1c2.png}
  \end{subfigure}
  \quad
  \begin{subfigure}[b]{0.45\textwidth}
    \centering
    \includegraphics[width=8cm]{../../../../Image/Book1/1-1c3.png}
  \end{subfigure}
\end{figure}

Cánh tay máy tự động nắm lại, sau đó úp xuống sàn và bắt đầu trườn đi, để lại một vệt ma sát rõ rệt trên nền nhà.

\gif{1-1d}{1}{179}

``Cánh tay của Roboco-senpai kìa\~{}!'' Mel và Fubuki chỉ tay hét lớn.

``Cánh tay của tôi\~{}!'' Roboco hoang mang kêu lên.

Hai cô nàng tóc vàng đồng loạt nhảy ra xa, còn Matsuri chỉ kịp kêu lên một tiếng nhỏ trước khi ngất xỉu ngay tại chỗ. Fubuki lập tức lao tới đỡ lấy Matsuri, hoảng hốt gọi tên cô: ``Matsuri-chan!?'' Trong lúc đó, nàng ``Đỏ'' và nàng Dơi lén liếc nhìn chiếc cánh tay máy vẫn đang trườn đi một cách kỳ dị, vẻ mặt đầy cảnh giác, như thể sợ nó sẽ tấn công bất cứ lúc nào.

Hoàn toàn trái ngược với sự kinh hoàng bao trùm cả căn phòng, Roboco chỉ cúi đầu nhìn xuống khớp vai trống hoác của mình, khẽ thở dài đầy lo lắng:

``Như thế này thì làm sao mình chơi game được nữa đây?''

``Đó là điều duy nhất chị đang lo lắng sao!?'' Mel ngơ ngác kêu lên, hai cánh dơi nhỏ sau lưng đập lia lịa vì sững sờ.

Haato rụt rè tiến lại gần, e dè nhìn vào phần vai bị đứt của tiền bối: ``A-nô\dots\ Chị có thấy đau không ạ?''

Roboco chớp chớp mắt, đáp bằng giọng bình thản đến lạ lùng: ``Không hề.''

Câu trả lời dứt khoát của cô nàng rô-bốt khiến Haato ngơ ngác, chỉ biết cười trừ vì không hiểu nổi cấu tạo cảm biến của một người máy công nghệ cao.

Lúc này, Fubuki đứng phắt dậy, vung tay tuyên bố dõng dạc: ``Dù sao đi nữa, chúng ta cũng phải bắt lấy cánh tay đó trước đã!''

Nghe vậy, mọi người nhanh chóng lấy lại tinh thần, cùng nàng Cáo trắng bắt đầu cuộc rượt đuổi cánh tay người máy đầy kỳ quặc. Tuy nhiên, mọi chuyện chẳng hề dễ dàng như họ tưởng. Vừa chạy được ba bước, cả ba cô gái bất ngờ vấp phải Matsuri, người vẫn còn nằm bất tỉnh nhân sự trên sàn, tạo thành một cú ngã chồng chéo đầy thảm hại.

\img{1-1e}

Cảnh tượng ba cô gái nằm đè lên nhau kèm theo những tiếng kêu oai oái khiến Roboco bất giác quay đầu lại nhìn. Vẻ mặt cô vẫn vô tư như thường lệ: ``Mấy đứa không sao chứ? Đừng để cánh tay của chị chạy xa quá nha.''

Haato nằm bẹp dưới cùng, cố sức ngóc đầu dậy gào lên trong bất lực: ``Ai\dots\ ai đó kéo Matsuri-chan ra giùm em với\dots!''

Fubuki vừa ôm đầu lồm cồm bò dậy, vừa chỉ biết méo mặt lẩm bẩm: ``Thật không thể tin nổi\dots\ cái cánh tay máy đó còn chạy nhanh hơn cả chúng ta nữa. Ái da\dots''

\ct

Sau khoảng mười phút hỗn loạn, các cô gái cuối cùng cũng dìu được Matsuri lên ghế sô-pha để nghỉ ngơi. Họ bắt đầu chia nhau đi lục soát khắp văn phòng. Fubuki lom khom kiểm tra từng ngóc ngách dưới gầm ghế, Mel quỳ hẳn xuống sàn để ngó vào gầm bàn, Haato lục lọi đống hộp các-tông chồng chất, còn Roboco thì điềm nhiên đứng tựa vào kệ tủ, mắt liên tục quét qua căn phòng. Nhưng dù có tìm kiếm mỏi mắt, cánh tay phải của cô nàng Người máy vẫn bặt vô âm tín.

``Có vẻ như nó đã trốn ra khỏi văn phòng rồi,'' Haato thở dài ngao ngán, đứng thẳng dậy phủi bụi trên quần áo.

Nghe vậy, Fubuki và Mel mệt mỏi thả mình xuống ghế xoay, mặt mày rầu rĩ như thể vừa để thua một trận đấu quan trọng.

Riêng nhân vật chính của vụ việc là Roboco thì vẫn bình thản ngồi bệt dưới sàn nhà. Cô nàng vừa ngân nga một giai điệu khe khẽ, vừa dùng bàn tay trái còn lại lướt điện thoại vô cùng chăm chú: ``Tay của mình chạy đi đâu rồi ta\~{}?''

Fubuki quay sang nhìn cô tiền bối, nửa bất lực nửa tò mò: ``Người bị rụng tay có vẻ là người ít lo lắng nhất ở đây nhỉ?''

Roboco nhún vai nhẹ, dửng dưng đáp: ``Thì tại chị lười đi tìm nó nữa rồi.''

Câu trả lời quá đỗi ngang ngược khiến Mel phải tròn mắt cảm thán: ``Hể\~{}? Sao chị có thể thản nhiên như vậy được chứ?''

``Thì đây có phải lần đầu chị bị rụng bộ phận đâu\~{} Trước đây tay chân chị cũng rơi rớt suốt mà\~{}'' Roboco thản nhiên trả lời.

Sự vô tư của cô nàng khiến cả ba người còn lại đồng loạt đổ mồ hôi hột. Haato chỉ biết xoa trán, tự hỏi lòng: ``Có khi nào mình cũng sẽ dần quen với những điều kỳ quặc thế này ở văn phòng không ta\dots?''

Nàng Ma cà rồng thở dài đầu hàng. Nhưng khi ánh mắt cô vô tình lướt qua chiếc thùng hàng lớn vẫn nằm yên vị cạnh lối vào, một tia sáng chợt lóe lên. Cô chỉ tay hỏi: ``Mà, không biết trong cái hộp to đùng kia chứa gì nhỉ?''

\img{1-1f}

``Trên hộp chẳng có ghi thông tin người gửi hay địa chỉ nhận gì cả\dots'' Roboco đáp, mắt vẫn không rời khỏi màn hình điện thoại.

``Ủa? Vậy rốt cuộc nó được gửi cho ai?'' Mel nghiêng đầu thắc mắc.

Đến lúc này, cô nàng rô-bốt mới giật mình nhận ra mình đã hoàn toàn bỏ qua khâu kiểm tra tem bưu phẩm. Cô nàng cười gượng, một tay gãi đầu hối lỗi: ``A-ha\dots\ Hình như chị quên kiểm tra mất tiêu rồi. Xin lỗi mọi người nha\~{}!''

Haato lắc đầu cười trừ, bước lại gần chiếc hộp để xem xét kỹ hơn. Sau một hồi săm soi, cô reo lên: ``Hộp này gửi đích danh cho chị đó, Roboco-senpai!''

``Ể\~{}? Thật á?'' Roboco chớp mắt đầy ngạc nhiên.

Không để mất thời gian, nàng Nữ sinh trung học nhanh chóng rạch băng keo và mở chiếc hộp ra. Ngay khi nắp hộp vừa hé mở, đôi mắt cô mở to vì kinh ngạc: ``Không lẽ\dots''

``\dots\ Đồ của Roboco-senpai thật sao?'' Fubuki cũng tò mò rướn người ngó vào.

Bên trong hộp là một cánh tay cơ khí hoàn toàn mới, sáng bóng và có thiết kế giống hệt cánh tay phải vừa trốn thoát của Roboco.

Cô nàng Người máy lập tức sáng mắt, lao tới cầm cánh tay lên và nở một nụ cười rạng rỡ: ``Đúng là cánh tay phải của chị rồi!''

Nói đoạn, cô nàng thản nhiên cắm thẳng cánh tay mới vào khớp vai, xoay xoay vài vòng rồi bấm nút tinh chỉnh. Khi nghe thấy tiếng khớp cơ khí khóa lại giòn giã, Roboco vung tay đầy phấn khích: ``Được rồi\dots\ ổn áp luôn!''

Mel đứng bên cạnh chứng kiến toàn bộ quá trình, há hốc mồm kinh ngạc: ``N-Nó\dots\ cứ thế gắn vào là xong thôi sao?!''

Trong khi Fubuki còn đang thẫn thờ vì sự tiện lợi của công nghệ ma giới thì cô chợt phát hiện dưới đáy hộp còn sót lại một mảnh giấy gấp gọn. Cô cúi xuống nhặt lên: ``Ơ, trong này còn có cả một lá thư nữa nè!''

``Đưa đây cho chị xem nào\~{}'' Roboco tò mò giật lấy bức thư từ tay Fubuki.

Cô nàng mở lá thư ra, nhanh chóng đọc lướt qua rồi cất giọng đều đều đọc to cho cả phòng cùng nghe:

\txtbx{
  Roboco-san à,

  Dạo này em khỏe chứ?

  Bọn anh hy vọng em sẽ vui khi nhận được món quà này. Mấy ngày này ở văn phòng có vui không? Nếu có, gửi lời chào của anh đến các thành viên trong công ty nhé.

  Dạo gần đây, anh để ý thấy em chơi game rất nhiều (đặc biệt là trò bắn súng), dẫn đến việc các bộ phận của em rất hay bị hỏng, đặc biệt là cánh tay phải.

  Vì biết rằng sẽ có một ngày nó bị rời ra, anh nghĩ em nên thay thế cánh tay đó bằng cái mà bọn anh đã gửi.

  Các bộ phận khác bọn anh sẽ tìm cách gửi sớm thôi, em đừng lo lắng quá nhé.

  Chúc em luôn khỏe mạnh!

  Thân gửi,

  \hfill\textbf{-Hikawa Kuon.}
}

Đọc xong lá thư, Roboco cẩn thận gấp gọn lại, nở một nụ cười ngượng ngùng: ``Hì hì\dots\ Hóa ra là linh kiện dự phòng được gửi trước à? Không ngờ Kuon-senpai lại có thể chế tạo được một cánh tay khớp hoàn hảo với chị đến vậy.''

Fubuki khoanh tay, bĩu môi hờn dỗi: ``Quà cáp gì chưa thấy, chỉ thấy làm bọn em một phen hú vía! Kuon-senpai thật là\dots!''

``Nhưng mà phải công nhận anh ấy chu đáo thật đấy!'' Haato vỗ vai Roboco cười khúc khích. ``Cứ như thể anh ấy có siêu năng lực biết trước tương lai cánh tay của chị sẽ rụng vậy!''

Roboco nhướng mày, khẽ lẩm bẩm tự bào chữa: ``Nhưng mà chị đâu có chơi game nhiều đến mức đó đâu ta\dots\ \textit{Sao em ấy lại biết được hay vậy nhỉ?}''

Bầu không khí căng thẳng ban nãy cuối cùng cũng được xua tan bằng những tiếng cười giòn giã của các cô gái khi nghĩ lại cảnh tượng hỗn loạn vừa rồi.

\ct

Sau khi dọn dẹp và sắp xếp lại đống lộn xộn trong phòng, mọi người bắt đầu quay trở lại với công việc của mình. Không khí trong văn phòng trở nên yên ả và dễ chịu hơn rất nhiều.

``Vậy ra cánh tay cũ của Roboco-senpai chỉ đơn giản là bị hỏng khớp nối thôi nhỉ,'' Haato vừa lau dọn mặt bàn vừa trò chuyện.

``Nhưng rốt cuộc tại sao Kuon-senpai lại phải gửi bưu phẩm tới đây làm gì nhỉ?'' Fubuki thắc mắc. ``Anh ấy hoàn toàn có thể tự tay mang đến văn phòng mà.''

Roboco bước lại gần, trên lưng cõng Matsuri vẫn còn đang ngủ say sưa. Mel đi bên cạnh, cười nói hùa vào: ``Chị cũng chịu thôi\~{} Nhưng mà ít nhất bây giờ Roboco-senpai có thể tiếp tục chơi game là tốt rồi!''

``Chính xác! Nhất là game Fortnite ấy, hôm nào rảnh cả lũ tụ tập bắn chung đi!'' Roboco vui vẻ đề xuất.

Haato khẽ nhíu mày, mặt hơi méo xẹo: ``N-Nhưng em có biết chơi bắn súng đâu\dots''

``Không sao đâu, chị gánh được tuốt! Chơi vài ván là nghiện ngay ấy mà!'' cô nàng Người máy vỗ ngực tự tin.

Fubuki mỉm cười nhìn mọi người, rồi chuẩn bị bước ra ngoài sảnh để pha trà. Thế nhưng mới đi được nửa đường, nàng Cáo bỗng khựng lại, khẽ nhíu mày đưa tay lên cằm suy nghĩ:

``Sao tự dưng mình cứ có cảm giác quên mất chuyện gì cực kỳ quan trọng ấy nhỉ\dots''
%==========================================TRUYỆN PHỤ=========================================
\omk

\begin{center}
  {\textit{*Các nhân vật nói bằng tiếng Etsunan.}}
\end{center}

\pov{Hikawa Kuon}

Tại xưởng cơ khí của gia đình Hikawa chúng tôi, hai bố con đang tất bật xử lý một vài đơn hàng tồn đọng của khách hàng. Xưởng cơ khí này là nơi gia đình tôi chuyên sửa chữa, bảo dưỡng và lắp ráp các thiết bị điện tử, máy móc cơ khí. Nó nằm cách nhà chúng tôi tận bốn phường, nghĩa là muốn đến đây làm việc, hai bố con phải lạch cạch đèo nhau bằng xe máy từ sáng sớm tinh mơ nếu không muốn chịu cảnh tắc đường kinh hoàng ở Kawachi.

Bình thường, xưởng này chỉ có bố tôi và người bác cả của tôi quán xuyến. Nhưng vì bác ấy có khu xưởng riêng chuyên làm đồ cơ khí nặng ở phía sau, nên chúng tôi chủ yếu hoạt động ở gian phía trước. Đợt này có đơn hàng cơ khí khẩn cấp mà chú đồng nghiệp của bố lại đột ngột xin nghỉ phép, nên bố bất đắc dĩ phải gọi tôi về phụ giúp một tay. Thú thật, tôi không phải là người quá am hiểu về kỹ thuật cơ khí, cũng chẳng thể vác nổi những món đồ sắt thép nặng trịch như bố, nhưng ít nhất tôi vẫn giúp được ông những việc vặt như đo đạc thông số, tính toán kích thước bản vẽ và đứng trông máy chạy.

Nghĩ đi nghĩ lại, chắc tôi phải sớm nghiên cứu lắp đặt một hệ thống tự động hóa cho cái xưởng này mới được. Chậc, làm thủ công thế này oải quá.

Đang lúc tôi lơ đãng nghĩ ngợi xem trưa nay hai bố con sẽ ăn gì thì đột nhiên--

*RẦM! RẦM! RẦM!*

``Ra ngay đây!'' tôi vội vàng hô lớn đáp lại. Tiếng đập cửa dồn dập vang lên từ phía cổng ngoài phá tan bầu không khí ồn ào của máy móc trong xưởng.

Tôi lật đật chạy ra ngoài kiểm tra thì bắt gặp một cánh tay máy đang nằm lăn lóc trên nền đất bám đầy cát bụi. Nhìn qua kiểu dáng công nghiệp và lớp sơn đặc trưng, tôi liền nhận ra ngay: Nó trông giống hệt như cánh tay của Roboco-san.

\dots Không, nói chính xác hơn thì đây chính là cánh tay của cô ấy chứ chẳng lệch đi đâu được! Nhìn kết cấu khớp nối đặc biệt này thì làm sao mà nhầm lẫn nổi! Nhưng mà sao nó lại tự bò đến tận đây được nhỉ?

À phải rồi. Hình như tôi chưa từng kể với ai là mình đã cài đặt một cơ chế khẩn cấp có tên ``tự động quay trở về nơi sản xuất'' cho các bộ phận cơ thể của cô ấy mỗi khi gặp sự cố nghiêm trọng thì phải?

Dù sao thì, tôi cũng nhặt cánh tay máy lên, tạm thời ngắt nguồn kích hoạt của nó rồi mang vào trong xưởng. Khoảng ba mươi phút sau, khi đã đến giờ nghỉ trưa, hai bố con tôi mới ngồi lại bên bàn làm việc và bắt đầu xem xét kỹ lưỡng ``vật thể lạ'' này.

\ct

Bố tôi chỉnh lại chiếc kính lão đã hơi trễ xuống mũi, dùng ngón tay thô ráp chỉ vào phần khớp gắn nối giữa vai và cánh tay của Roboco rồi khẽ lắc đầu thở dài: ``Bố chịu rồi\dots\ Cái này trông có vẻ không sửa nổi đâu con.''

``Sao thế ạ bố?'' tôi tò mò ghé đầu vào nhìn.

``Cái chốt chịu lực chính này bị gỉ sét ăn mòn hết cả rồi. Nếu muốn thay nó thì bắt buộc phải tháo rời toàn bộ cụm chuyển động này ra, mà đoạn này họ lại gia cố bằng đinh tán cố định từ lúc đúc khuôn rồi. Tháo ra là hỏng cấu trúc bên trong ngay, không cách nào phục hồi nguyên vẹn được,'' bố tôi chậm rãi giải thích.

\chrint

Người đang đưa ra chẩn đoán đầy chuyên nghiệp đó là bố tôi - Hikawa Ryujin (\ruby{火}{ひ}\ruby{川}{かわ}\ruby{龍}{りゅ}\ruby{人}{じん}, Hỏa Xuyên Long Nhân). Ông là một thợ cơ khí lành nghề đã ngoài 55 tuổi với hơn 40 năm kinh nghiệm trong ngành, tính từ lúc ông còn là một cậu học sinh trung học phụ việc ở xưởng. Thực chất, cái nghề lấm lem dầu mỡ này là do đời ông nội tôi truyền lại, và bản thân bố tôi từ nhỏ cũng đã có niềm đam mê đặc biệt với các loại máy móc cơ khí nên đã quyết định nối nghiệp gia đình.

Hai bố con tôi từ lâu đã vô cùng thân thiết, thường được mọi người xung quanh ví như hai ``bạn vong niên'' hơn là quan hệ bố con thông thường. Có lẽ vì tôi và em gái sống cùng bố nhiều hơn với mẹ. Hồi tôi lên mười tuổi, bố mẹ từng có một khoảng thời gian dài ly thân - nhấn mạnh là ly thân do bất đồng quan điểm gia đình chứ không phải ly hôn. Có lẽ cũng nhờ khoảng thời gian tự lập cùng bố mà hai bố con hiểu nhau hơn. Ngay cả đôi mắt với con ngươi màu đỏ đặc trưng của tôi cũng là được di truyền hoàn toàn từ ông.

Trang phục thường ngày của bố lúc nào cũng bị mẹ chê là luộm thuộm: Một chiếc áo phông mỏng bạc màu đi kèm quần jean và đôi dép lê giản dị vào những ngày hè oi bức; còn khi trời chuyển lạnh thì sẽ là áo len, quần nỉ dày và giày thể thao cũ. Nói một cách đơn giản, đó là bộ trang phục thoải mái nhất để ông có thể dễ dàng chui gầm xe hay sửa máy bất cứ lúc nào.

\ct

Nhìn vào biểu cảm của bố, tôi cũng đã tự hiểu rõ vấn đề. Muốn cứu vãn cánh tay của cô tiền bối Robo-PON thì chỉ có cách đập đi xây lại hoàn toàn, chứ không thể sửa chữa bảo dưỡng thông thường. Mà nếu có cố sửa thì chi phí linh kiện thay thế và thời gian bỏ ra cũng quá lớn. Nói cách khác, việc này hoàn toàn không thực tế.

Bố tôi thở dài, châm một điếu thuốc lá, trầm ngâm nhìn cánh tay máy đang nằm im lìm trên bàn: ``Xưởng mình lại đang dồn dập đơn hàng của khách cần bàn giao gấp trong hôm nay. Bố không nghĩ chúng ta có đủ thời gian để cố đấm ăn xôi sửa cái này đâu.''

``Tức là bây giờ\dots\ bố con mình phải\dots'' tôi ngập ngừng.

``Tiêu hủy nó thôi con ạ. Cố sửa cũng chỉ phí công vô ích,'' bố tôi dứt khoát đưa ra kết luận.

``Nhưng\dots\ chúng ta thật sự không còn cách nào giúp nó sao bố? Nó đã tự bò một chặng đường dài về tận đây mà\dots'' tôi nhìn cánh tay máy, trong lòng dâng lên một chút gợn sóng tiếc nuối.

``Thôi con ạ, cuộc sống là thế. Đôi khi có những thứ hư hỏng nặng đến mức dù ta cố gắng đến mấy cũng không thể khôi phục lại như ban đầu được nữa,'' bố nhẹ nhàng đặt tay lên vai tôi trấn an.

Đúng lúc này, cánh tay của Roboco dường như vừa tự kích hoạt lại nguồn năng lượng dự phòng cuối cùng. Nó khẽ cử động các ngón tay, nhưng có vẻ như đã nghe hiểu được cuộc trò chuyện của hai bố con tôi. Biết trước số phận của mình đã đến lúc phải chấm dứt, cánh tay chỉ khẽ run lên rồi buông thõng xuống tỏ vẻ cam chịu, dù không thể thốt ra bất kỳ lời nào.

Tôi nhìn cánh tay ấy, hít một hơi thật sâu rồi quả quyết nói với bố: ``Vâng. Nhưng mà bố biết không? Con sẽ tiễn đưa cánh tay này theo cách mà nó mong muốn.''

\ct

Tôi mang cánh tay máy ra khu vực lò luyện kim ở phía sau xưởng. Đặt nó nằm gọn lên một chiếc dây xích sắt treo sẵn, tôi khẽ gật đầu ra hiệu rồi từ từ điều khiển hệ thống ròng rọc hạ dần dây xích xuống miệng lò - nơi dòng dung nham kim loại nóng chảy đang cuồn cuộn đỏ rực phía dưới.

Cánh tay từ từ chìm dần vào làn hơi nóng hầm hập. Và rồi, ngay trong khoảnh khắc chuẩn bị tan chảy hoàn toàn, nó bất ngờ dùng chút động lực cuối cùng, nắm chặt các ngón tay lại rồi từ từ dựng đứng ngón tay cái lên hướng về phía tôi - một cái thumbs-up đầy kiên định và quả cảm trước khi chìm hẳn\footnote{Phân cảnh tiễn biệt huyền thoại trong bộ phim Terminator 2, nơi kẻ hủy diệt T-800 chấp nhận chìm vào bể gang lỏng để bảo vệ tương lai nhân loại.}.

Lòng dâng trào cảm xúc, tôi đứng thẳng người, đưa tay lên thái dương chào theo kiểu nhà binh để tiễn đưa người bạn cơ khí trung thành lần cuối, cho đến khi ngón tay cái của nó biến mất hoàn toàn trong ánh sáng chói lòa của kim loại nóng chảy.

Đứng trước lò dung nham cuồn cuộn, tôi lặng người nhìn tàn tro bay lơ lướng, cảm giác hụt hẫng và mất mát dâng tràn trong lồng ngực. Dẫu biết rằng đây chỉ là một bộ phận cơ khí bị hỏng không thể phục hồi, tôi vẫn cảm thấy nặng lòng khi chính tay mình phải tiễn đưa nó. Nhưng tôi hiểu đây là cách giải thoát duy nhất cho nó. Nó đã từng hoàn thành xuất sắc nhiệm vụ của một phần cơ thể hoàn chỉnh, và giờ đây nó đã trở về với cát bụi. Rồi đây, Roboco-san chắc chắn sẽ có những bộ phận mới tốt hơn, hoàn thiện hơn để thay thế.

Bố tôi đứng bên cạnh từ nãy đến giờ, không hề lên tiếng cắt ngang, chỉ lặng lẽ quan sát thái độ của con trai. Tôi biết ông cũng đang có những suy ngẫm riêng của mình.

``Con thấy không, Kuon?'' giọng bố trầm xuống, vang vọng trong không gian gian xưởng rộng lớn. ``Không phải thứ gì trong cuộc đời này cũng có thể sửa chữa được. Đôi khi, học cách chấp nhận sự thật và buông tay để chúng ra đi mới là lựa chọn tốt nhất.''

Ông nhìn tôi với ánh mắt trìu mến của một người đi trước đã trải qua bao sóng gió: ``Cuộc sống là một chuỗi của những điểm kết thúc và những khởi đầu mới. Việc của chúng ta là làm tốt nhất có thể ở hiện tại, rồi tiếp tục vững bước về phía trước.''

Tôi khẽ gật đầu, thầm hiểu rằng lời bố nói không chỉ đơn thuần là dành cho cánh tay máy của Roboco. Ông đang nói về triết lý nhân sinh, về cách đối mặt với những khó khăn trong đời và cách học cách từ bỏ những điều không thể thay đổi. Tôi lặng lẽ khắc ghi bài học sâu sắc ấy vào lòng, tự nhủ bản thân phải mạnh mẽ hơn để đối mặt với những biến cố trong tương lai.

Cuối cùng, hai bố con chúng tôi cùng nhau rời khỏi khu vực lò luyện kim. Cả hai bước đi trong im lặng, mỗi người mang một suy nghĩ riêng, nhưng trong lòng đều thấu hiểu rằng sự mất mát này không phải là dấu chấm hết, mà chỉ là một phần tất yếu của hành trình trưởng thành.
\end{document}